\documentclass[a4paper, 12pt, article]{memoir}

\setlength{\footmarkwidth}{0em}
\setlength{\footmarksep}{0em}
\setlength{\thanksmarkwidth}{0em}
\setlength{\thanksmarksep}{0em}

\usepackage{fontspec}
\usepackage{amsmath}
\usepackage[tracking=true, nopatch=footnote]{microtype}
\usepackage{scalefnt}
\usepackage{csquotes}
\usepackage[main=english]{babel}

\usepackage{graphicx}
\usepackage{subcaption}
\usepackage{tikz} %for all basic options
\usepackage{tikz-qtree} %for simple tree syntax
\usepgflibrary{arrows} %for arrow endings
\usetikzlibrary{positioning,shapes.multipart} %for structured nodes
\usetikzlibrary{tikzmark}
\usepackage{tikz-qtree}
\usepackage{tree-dvips}
\usepackage{booktabs}
\usepackage{longtable}
\usepackage{array}
\usepackage{multirow}
\usepackage{wrapfig}
\usepackage{float}
\usepackage{colortbl}
\usepackage{pdflscape}
\usepackage{tabu}
\usepackage{threeparttable}
\usepackage{threeparttablex}
\usepackage[normalem]{ulem}
\usepackage{makecell}
\usepackage{xcolor}

\setmainfont{Gentium Plus}
\setmonofont[Scale=0.8]{Noto Sans}

\babelprovide[import, onchar=ids fonts letters]{ancientgreek}
% \babelfont[ancientgreek]{rm}{Brill}
\babeltags{grc = ancientgreek}

\usepackage{hyperref}
\usepackage{linguex}
\setlength{\Exlabelsep}{0.5em}
\setlength{\SubExleftmargin}{1.5em}
\usepackage{paralist}
\usepackage{hyphenat}
\hyphenation{
  Κυ-ρι-α-κί-δη
  Θεσ-σα-λο-νί-κη
}

\usepackage[normalem]{ulem}
\RequirePackage{xcolor}
\definecolor{green}{RGB}{16,87,87} % rgb(16,87,87)
\definecolor{red}{RGB}{193, 11, 105} % rgb(193, 11, 105)
\definecolor{yellow}{RGB}{218,222,104} % rgb(218,222,104)
\definecolor{pink}{RGB}{243,179,145} % rgb(243,179,145)
\definecolor{blue}{RGB}{161,184,206} % rgb(161,184,206)

\RequirePackage{hyperref}%
\hypersetup{%
	colorlinks=true, % false: boxed links; true: colored links
	linkcolor=green,  % color of internal links
	citecolor=green,  % color of links to bibliography
	filecolor=green,  % color of file links
	urlcolor=green,
}

\usepackage[
  backend=biber,
  uniquename=init,
  giveninits,
  dashed=true,
  style=authoryear-comp
  ]{biblatex}
\addbibresource{./biblio.bib}


\title{Rethinking case attraction on Ancient Greek infinitive clauses}
\author{
  Caio B. A. Geraldes\thanks{Fundação de Amparo à Pesquisa do Estado de São
  Paulo.}
	\\\smallskip
	<\href{mailto:caio.geraldes@usp.br}{\texttt{caio.geraldes@usp.br}}>
}
\date{June 25, 2025}

\begin{document}

\maketitle

\chapter{Introduction} % (fold)

\paragraph{Good morning.} 
Nominal predicates and adjuncts modifying subjects of complement infinitive
clauses in Ancient Greek (AG) usually agree in case with their domain local
subject, which has morphological accusative case assigned in the
\emph{Accusativus cum Infinitivo} (\emph{AcI}) construction.
When the subject of the infinitive is not overtly expressed but the main clause
has a dative or genitive object which shares the reference of the infinitive
subject, these predicates or adjuncts may either occur in the accusative or ``be
attracted'' to the case of the main clause's object.

\ex.\label{gloss:base}
\a.\label{gloss:noattr}συμβουλεύει \textbf{τῷ Ξενοφῶντι ἐλθόντα} {εἰς Δελφοὺς}
ἀνακοινῶσαι {τῷ θεῷ} {περὶ τῆς πορείας}.\\
He advices Xenophon to go to Delphi and ask the god about the travel.
(Xen.\ \emph{Anab.} 3.1.5)
\b.\label{gloss:attr}ἀφῆκε \textbf{μοι ἐλθόντι} {πρὸς ὑμᾶς} λέγειν τἀληθῆ.\\
He allowed me to go and speak the truth in front of you.
(Xen.\ \emph{Hell.} 6.1.13)


The literature on the subject, on the more philological side, from
Port\hyp{}Royal (\cite{Lancelot1658,Lancelot1709}) up to the Cambridge
Grammar~(\cite{Boas2019}), have focused on listing the contexts in which
attraction may occur, whereas the transformational and minimalist approaches
of~\textcite{Lakoff1970}, \textcite{Andrews1971}, \textcite{Quicoli1982},
\textcite{Spyropoulos2005}, \textcite{Tantalou2003},
\textcite{Sevdali2006,Sevdali2013} have mainly occupied themselves with
proposing the structure that would allow the phenomenon to occur.
What factors influence the decision between either resolution
in~\ref{gloss:base} remains, however, largely unknown.

\paragraph{[Self-deprecatingly]} I myself have, in the last edition of this
same event, tried to present the distributional patterns that make the
attraction,~\ref{gloss:attr}, more likely, showing that there is strong
correlation between attraction and the classes of the matrix and infinitive
verbs and the relative distance and position of the controller and target of
agreement.

Furthermore, I argued that these factors were \emph{proxies} to unobserved
values, namely semantic and pragmatic features of the discourse setting, and
that even in the absence of seen contrasting factors, precisely the case in the
pair above~\ref{gloss:base}, the unseen discourse factors would be still causing
one resolution to be at the very least more likely.

\paragraph{[Introductingly]}
I wasn't, however, quite convinced about neither the methods, the data, and,
consequently, the quality of the results at which I had arrived.
The main reason for this distrust is the fact that there could be interaction
between the factors analysed for which the methods employed were not well designed
to account, at least not with the amount of data available.
A secondary motive was the classic adage ``correlation is not causation''.

To solve these shortcomings, I have collected more data from a larger
\emph{corpus} but more importantly swapped the quantitative framework for one
capable of dealing with multivariate interactions, namely from ``out of the
shelf'' statistical tests to generalized linear models, and, I believe more
importantly, to for one capable of expressing the causality behind the
statistical correlations.

For this to work, however, we must take some steps back and build a causal
model for what is the linguistic nature of case attraction and how properties of
the sentence and its constituents might favor its occurrence.
The following is a walk-through the construction of this causal model and its
results.

\chapter{The data} % (fold)

Before stepping into the framework and analysis proper, a brief word on the
empirical data supporting this research and any conclusion coming from it.
The dataset is composed of 343 sentences drawn from the Attic literary
\emph{corpus} and Herodotus.
The sentences were extracted from the Diorisis Corpus
(\cite{TheDiorisisAncientGreekCorpus,diorisis2021}) with help of the project's
querying tool, Diorisis Search (\cite{DiorisisSearch}).
Although comprehensive, the selection is by no means exhaustive.

Each sentence has been annotated for values of interest, including authorship,
literary genre, dialect as well as class of matrix and infinitive verb, class of
controller and target, distance between constituents, etc.
Parts of the annotation processes were conducted programmatically and later
checked manually.

Every model used in this work was processed using R (\cite{R}), the packages
\texttt{tidyverse} (\cite{tidyverse}), \texttt{dagitty} (\cite{dagitty}),
\texttt{ggdag} (\cite{ggdag}), and \texttt{rethinking} (\cite{McElreath2020}),
as well as the software Stan~(\cite{StanReferenceManual,StanUserGuide}).

% chapter The data (end)

\chapter{Nature of attraction} % (fold)

In the most abstract way possible, case attraction in infinitive clauses,
exemplified by~\ref{gloss:attr}, is when two constituents, a controller and a
target, produced in different domains of a sentence, one of the domains
encompassing the other, covary in an specific feature, namely $[Case]$.
This satisfies the definition of agreement in Corbett's typological assessment
of the phenomenon across languages (\citeyear[pp.\ 4--7]{Corbett2006}).

The lack of covariation of the $[Case]$ feature, seen in~\ref{gloss:noattr}, is
generally understood as result of the application of \textbf{default} case
concord rules of Ancient Greek.
This assumes that the subject of infinitive clauses in Greek is always marked as
accusative, even when not overtly expressed, and that convert constituents are
also marked for $[Case]$ and its case is also available for nominal concord.

The facts that the target has a choice of controller (either the overt oblique or
de convert accusative) for agreement, that $[Case]$ is a non\hyp{}lexical
features, that the agreement is optional, and that the domain is non\hyp{}local
favours the interpretation that the agreement type of case attraction is less
canonical than the general agreement type of Ancient Greek, following the
criteria for canonical agreement in the same work of \textcite[pp.\
8ff.]{Corbett2006}.

If this assessment is sound, it is expected that attraction and
non\hyp{}attraction behave as competing agreement resolutions, the former being
less canonical and thus more likely to be conditioned.
Conditions concerning the semantic and pragmatic properties of the controller,
its precedence in relation to the target and the class of predication (i.e.:
class of the domain) have been identified by Corbett.
Although properties of the target were not shown to be favouring any resolution
in agreement mismatches in Corbett's survey, case attraction might be one
example of it, and I will hypothesize that properties of targets are competing
causes of non\hyp{}canonical resolutions.
A model of agreement in this framework may be represented as a direct acyclic
graph, in which \textbf{A}ttraction ($\simeq$ Agreement) is caused by features
of the \textbf{C}ontroller, \textbf{D}omain and \textbf{T}arget:

\begin{center}
	\includegraphics[width=0.40\textwidth]{../analysis/figs/dag_non_canonical_agreement.png}
\end{center}

These Controller, Domain and Target nodes represent clusters of linguistic
factors, some of which we may observe directly (part of speech of the target),
some might be implied (animacy) and some may remain unobserved (discourse
function, as Ancient Greek does not overtly marks focuses or topics in the
morphosyntax, at least not dependably).
When an observed factor is included in the model, we split the cluster into
two nodes: the one of observed factor and that containing rest of unobserved features.

The three factors previously shown to correlate to case attraction represent
each of these nodes:
\begin{compactitem}
	\item Copular infinitives imply a different domain of agreement, as they
		select different predication classes;
	\item Main verbs with semantic modality imply a different type of controller,
		as they have higher subjecthood;
	\item Word order and distance imply different status for the target, as
		targets more to the left of the sentence are both structurally closer to the
		controller and have a higher communicative function.
\end{compactitem}

% chapter Nature of attraction (end)

\chapter{The domain of attraction} % (fold)

Firstly, \textcite{KG} and others have associated case attraction with sentences
in which the infinitival clause is a copular infinitive, einai or genesthai.
Should the class of the infinitive favour either resolution of agreement?
Probably not directly, as there is no reason linguistically it would.
Copular infinitives, however, have a double effect on the agreement ``causes'':
\begin{compactenum}
\item they filter the possible relations between controller and target, as those are
now more likely in primary predication, rather than secondary or appositive
relation;
\item this then delimits the option of target, making it more likely that they are nouns or adjectives, rather than participles.
\end{compactenum}

The these are known to affect agreement resolutions, and a couple of Hierarchies
have been proposed to predict when the less canonical resolution (also called
semantic agreement) is more likely to take place. The Agreement Hierarchy
(\cite{Corbett1979}) and Predicate Hierarchy (\cite{Comrie1975}) have also been
combined into the Agreement and Predicate Hierarchy shown
in~\ref{hierar:agreementpredicate2}:

\ex.\label{hierar:agreementpredicate2} Agreement and Predicate Hierarchies~(\cite[233]{Corbett2006}):\\
\begin{tikzpicture}[baseline=0pt]
    \node (attr) at (0, 0) {Attributive};
    \node (pred) at (2.5, 0) {Predicative};
    \node (predv) at (2.9, -0.4) {Verb};
    \node (predp) at (3.5, 1) {Participle};
    \node (preda) at (5, 2) {Adjective};
    \node (predn) at (7, 2.5) {Noun};
    \node (relp) at (5.6, 0) {Relative Pronoun};
    \node (perp) at (9.1, 0) {Personal Pronoun};
    \node (less) at (0, -1) {-};
    \node (plus) at (9.1, -1) {+};
    \draw[->, thick] (attr) to (pred);
    \draw[->, thick] (pred) to (relp);
    \draw[->, thick] (relp) to (perp);
    \draw[->, thick] (pred) to (predp);
    \draw[->, thick] (predp) to (preda);
    \draw[->, thick] (preda) to (predn);
    \draw[->] (less) to (plus);
    \draw[->] (plus) to (less);
    \node at (4.5, -1.5) {Likelihood of agreement with semantic justification};
\end{tikzpicture}


Adding this to our causal model, we may say that \textbf{A}ttraction is caused
both by whether or not the relation between the controller and target is
copular, and by the Part of Speech of the target, which is itself ``caused'' by
the type of relation (copulas select more nominal predicates on average), and
other unobserved properties of the \textbf{D}omain may also cause Attraction.

By applying the do-calculus~(\cite{Pearl2009}), we learn that:
the direct effect of the part of speech class requires us to
adjust the statistic model by $C$, so that we block the non\hyp{}causal path $P
\leftarrow C  \rightarrow A$.

\begin{center}
\includegraphics[width=0.55\linewidth]{../analysis/figs/dag_a.png}
\end{center}



\chapter{The controller of attraction} % (fold)

Another association is that between case attraction and impersonal verbs.
I have, however, shown that this association does not exist with every single
impersonal verb, and that the verb δοκέω almost never heads a sentence with case
attraction.
The impersonal verbs which are actually associated with case attraction are
those carrying modal semantics, such as ἐξαρκεῖ and ἔξεστί.
Is it likely that the nature of the matrix verb directly affects the choice of
agreement? I don't think so, but it may indirectly do, by defining which
features the controller of agreement, the oblique object, has.
These verbs have very little semantic information and the theta role and
grammatical function of their oblique object seems to depend on the infinitive
verb.
We may say that σοι in~\Next, for example, is the subject of the sentence
(a quirk subject of sorts) and has the thematic role of agent, much in contrast
if the sentence was constructed with a verb of order or advice like συμβουλέυω,
to which the oblique represents a recipient of the action, and if the sentence
is constructed by δοκεῖ, to which the oblique represents an experiencer, ``the
one who believes X to be better to do''.

\ex. ἀλλ᾽ ἐξαρκῇ σοι τὰ ἔργα αὐτῶν ὁρῶντι σέβεσθαι καὶ τιμᾶν τοὺς θεούς.\\
But it is enough for you to look at these deeds and respect the gods. 
(Xen.\ \emph{Mem.} 4.3.13)

This, however, only holds if the infinitive verb is not copula, otherwise, the
thematic role is also that of an experiencer, meaning that \textbf{Th}ematic
Role is a sister node of \textbf{P}art \textbf{o}f \textbf{S}peech, as it is
``caused' by \textbf{C}opula.

Adding to our previous DAG, the model will look like:

\begin{center}
\includegraphics[width=0.55\linewidth]{../analysis/figs/dag_b.png}
\end{center}


\chapter{Intermezzo: Authors and dialects} % (fold)

I have argued previously that, in the corpus of Herodotus, Plato and Xenophon,
the verbs of modal semantics were preferred by the authors when writing in the
philosophical dialogue genre, and that this had implications if we were to estimate
the effect of authorship or dialect on case attraction.
If modality verbs affect the likelihood of attraction by means of selecting the
thematic role of the controller, the author's idiosyncrasies
in the use of attraction or their dialect base likelihood of using attraction
would be biased in the data by the effects of the selection of the main verb in
sentences where attraction may appear.

We must thus add \textbf{Auth}orship, \textbf{Genre} and \textbf{Dia}lect to our
causal model.

\begin{center}
\includegraphics[width=0.55\linewidth]{../analysis/figs/dag_c.png}
\end{center}


These last couple of additions have the following consequences:
\begin{compactitem}
\item to estimate the direct effect of \textbf{Th}ematic role, we must control
	for \textbf{Copula} and either by \textbf{Genre} or \textbf{Auth}orship;
\item to estimate the direct effect of \textbf{Cop}ula, we must control for
	\textbf{P}art \textbf{o}f \textbf{S}peech, \textbf{Th}ematic role, and either
	by \textbf{Genre} or \textbf{Auth}orship;
\item to estimate the direct effect of \textbf{Dia}lect, we must control for
	\textbf{Auth}orship;
\item to estimate the direct effect of \textbf{Auth}orship, we must control for
	\textbf{Dia}lect and by either \textbf{Genre} or \textbf{Cop}ula and
	\textbf{Th}ematic role.
\end{compactitem}

\clearpage

% chapter Intermezzo: Authors and dialects (end)


% chapter The controller of attraction (end)

\chapter{The target of attraction} % (fold)

Finally, it has been long been proposed that case agreement is more likely when
controller and target are very close to each other, and the likelihood
diminishes as it grows.
Furthermore, precedence of the target in relation to the controller
seems to demand case attraction, such as in~\Next.

\ex.ὡς \textbf{ἀγαθοῖς} τε \textbf{ὑμῖν} προσήκει εἶναι.\\
That you may be good people. (Xen.\ \emph{Anab.} 3.2.11)


Word order is highly associated in Ancient Greek with an important
condition for non\hyp{}canonical agreement: communicative
function.~\footnote{\textcite{Dik1995,Dik2007,Matic2003,Allan2012}}.
Although associated with, however, word order does not completely determines
communicative function of the target: it is not enough to guarantee that a
certain constituent is a topic or focus.
In other words, Word Order would only affect case attraction indirectly, by
`causing' the communicative function of the target, function which is, however,
unobservable.
This means that the effect of word order over case attraction is statistically
impossible to infer without a way to measure the communicative function of the
target.

The effect of distance, on the other hand, is measurable, so long as we control
for its effect on word order: distance changes word order, so we must block the
non\hyp{}causal path that flows from distance to attraction going through word
order and communicative function.
In the DAG\@:

\begin{center}
\includegraphics[width=0.60\linewidth]{../analysis/figs/dag_d.png}
\end{center}
\clearpage


% chapter The target of attraction (end)

\chapter{Analysis} % (fold)

Fortunately, we don't need many different models to estimate the effects of many
of the factors of interest.
By applying the do-calculus, so long as we control for Word Order, we can
estimate the direct effects of:
\begin{inparaenum}
	\item Part of Speech of the target;
	\item Copular infinitive (i.e.:\ class of predication);
	\item Thematic role of the oblique object;
	\item Dialect;
	\item Authorship; and
	\item Distance.
\end{inparaenum}
The following generalized linear model allows us to do so:

\begin{align*}
	\text{A} &\sim \text{Bernoulli}(p_i)\\
	\text{logit}(p_i) &= \alpha + 
	\beta_{\text{Cop}}\text{Cop}_i + 
	\beta_{\text{Dist}}\text{Dist}_i + \beta_{\text{WO}} \text{Prec}_i \\&+
	\beta_{\text{PoS}}\sum^{PoS_{i}-1}_{j=0}{\delta_{j_{\text{PoS}}}} + \beta_{\text{Th}}\sum^{Th_{i}-1}_{j=0}{\delta_{j_{\text{Th}}}} \\&+
		\beta_{\text{Auth}_{i}} + \beta_{\text{Dia}_i}\\
	\alpha &\sim \text{Normal}(0, 1)\\
	\beta_{\text{Cop}}, \beta_{\text{Dist}}, \beta_{\text{Prec}} &\sim \text{Normal}(0, 1)\\
	\beta_{\text{PoS}}, \beta_{\text{Th}} &\sim \text{Normal}(0, 0.5)\\
	\delta_{\text{PoS}}, \delta_{\text{Th}} &\sim \text{Dirichlet}(2)\\
	\beta_{\text{Auth}} &= z_{\text{Auth}}*\sigma_{\text{Auth}}\\
	\beta_{\text{Dia}} &= z_{\text{Dia}}*\sigma_{\text{Dia}}\\
	z_{\text{Auth}}, z_{\text{Dia}} &\sim \text{Normal}(0,1)\\
	\sigma_{\text{Auth}}, \sigma_{\text{Dia}} &\sim \text{Exponential}(1)
\end{align*}

A brief note on priors:
Those are fairly flat, as there is no reason to believe we have weak or
strong relationships in agreement phenomena.
The effects of Part of Speech and Thematic Role (the scary ones with sums) are
modeled as Ordered Categorical Predictors using the Dirichlet distribution.
This is to represent the generally agreed idea that typological hierarchies have
``cutoff'' points in between the nodes and that their effect increases
``monotonically'' as we traverse between the nodes.
The first node of each hierarchy will have the base effect $\beta_{X}$, and the
following nodes will ``accumulate'' effect. Using Part of speech as an example,
$\beta_{\text{PoS}}$ is the general effect of the part of speech over the
likelihood of Attraction, and $\delta$ is a vector with values denoting how much
going up from Participle to Adjective to Noun increases $\beta_{\text{PoS}}$.
Participles do not change the effect, Adjectives add $\delta_1$ and Nouns adds
both the effects $\delta_1 + \delta_2$.

The effects Authorship and Dialect have non\hyp{}centered priors to preform
partial pooling, so that we are able to estimate how much variation there is
between each of their values.

\section{Copula and Distance} % (fold)

The most generally associated factors with attraction, namely the infinitive
being a copula and the distance between controller and target show the expected
effects:

\begin{center}
	
\begin{tabular}[t]{lrrrrr}
\toprule
Effect & mean & sd & 5.5\% & 94.5\% & rhat\\
\midrule
$\beta_{\text{Copula}}$ & 1.81 & 0.47 & 1.07 & 2.57 & 1.00\\
$\beta_{\text{Dist}}$ & -0.27 & 0.06 & -0.37 & -0.18 & 1.01\\
\bottomrule
\end{tabular}

\end{center}

\begin{figure}[!h]
\centering
\begin{minipage}{.48\textwidth}
  \centering
  \includegraphics[width=\linewidth]{../analysis/figs/posterior_b_copula.png}
	\captionof{figure}{Posterior distribution of $\beta_\text{Cop}$ with the 89\%
	Highest Posterior Density Interval.}
\end{minipage}%
\hfill
\begin{minipage}{.48\textwidth}
  \centering
  \includegraphics[width=\linewidth]{../analysis/figs/posterior_b_dist.png}
	\captionof{figure}{Posterior distribution of $\beta_\text{Dist}$ with the 89\%
	Highest Posterior Density Interval.}
\end{minipage}
\end{figure}

The estimate effect for Distance is quite precise, and, to offer an intuition of
what these values mean, the model predicts that by the time the controller and
target are 8 words apart, on average, the fact that the infinitive is a copula
is no longer relevant, i.e., the attraction is no longer a more likely outcome.
If our causal model is sound, this is an important result for the typology of
agreement in general, as the evidence from other languages predict that higher
distances between controller and target would favor non\hyp{}canonical agreement
resolutions (e.g.:\ modifiers of committee nouns in English are more likely to
stand in the plural when far away from their controller, as shown by
\textcite[98]{Levin2001}).\footnote{Other examples in
\textcite[235]{Corbett2006}}.

% section Copula and Distance (end)

\section{Part of Speech and the Predicate Hierarchy} % (fold)

The estimated effects are described in the following table and plot.

\begin{center}
	
\begin{tabular}[t]{lrrrrr}
\toprule
Effect & mean & sd & 5.5\% & 94.5\% & rhat\\
\midrule
$\beta_{\text{PoS}}$ & 0.76 & 0.51 & 0.17 & 1.71 & 1\\
$\delta_{\text{Adj}}$ & 0.48 & 0.18 & 0.19 & 0.78 & 1\\
$\delta_{\text{Noun}}$ & 0.52 & 0.18 & 0.22 & 0.81 & 1\\
\bottomrule
\end{tabular}

	\bigskip
	
\begin{tabular}[t]{llrrrrr}
\toprule
PoS & Cummulative effect & mean & sd & 5.5\% & 94.5\% & rhat\\
\midrule
Participle & $\beta_{\text{PoS}} * 0$ & 0.00 & 0.00 & 0.00 & 0.00 & --\\
Adjective & $\beta_{\text{PoS}} * \delta_{\text{Adj}}$ & 0.35 & 0.26 & 0.06 & 0.86 & 1\\
Noun & $\beta_{\text{PoS}} * (\delta_{\text{Adj}} + \delta_{\text{Noun}})$ & 0.76 & 0.51 & 0.17 & 1.71 & 1\\
\bottomrule
\end{tabular}

\end{center}


\begin{center}
	\includegraphics[width=0.70\linewidth]{../analysis/figs/model_domain_banddelta.png}
\end{center}

To put it into words, the data favours the idea that the Agreement and
Predication Hierarchy holds for case attraction in Ancient Greek case
attraction: as the target becomes more ``nominal'', the likelihood of attraction
increases from basically neutral with participles up to highly positive.
There is still a fair amount of uncertainty on how strong the effect is,
however, and if the major cutoff point is between Participle and Adjective or
between Adjective and Noun.


\section{Thematic Role of the Controller} % (fold)

The estimated effects are described in the following table and plot.

\begin{center}
	
\begin{tabular}[t]{lrrrrr}
\toprule
Effect & mean & sd & 5.5\% & 94.5\% & rhat\\
\midrule
$\beta_{\text{Theta}}$ & 2.73 & 0.60 & 1.88 & 3.76 & 1\\
$\delta_{\text{Exp}}$ & 0.38 & 0.12 & 0.19 & 0.57 & 1\\
$\delta_{\text{Agent}}$ & 0.62 & 0.12 & 0.43 & 0.81 & 1\\
\bottomrule
\end{tabular}

	\bigskip
	
\begin{tabular}[t]{llrrrrr}
\toprule
Theta & Cummulative effect & mean & sd & 5.5\% & 94.5\% & rhat\\
\midrule
Recipient & $\beta_{\text{Th}} * 0$ & 0.00 & 0.00 & 0.00 & 0.00 & --\\
Experiencer & $\beta_{\text{Th}} * \delta_{\text{Exp}}$ & 1.04 & 0.41 & 0.44 & 1.73 & 1\\
Agent & $\beta_{\text{Th}} * (\delta_{\text{Exp}} + \delta_{\text{Agent}})$ & 2.73 & 0.60 & 1.88 & 3.76 & 1\\
\bottomrule
\end{tabular}

\end{center}

\begin{center}
	\includegraphics[width=0.75\linewidth]{../analysis/figs/model_theta_banddelta.png}
\end{center}

Again, the results confirm the hypothesis that attraction is more likely to
occur when the oblique object of the matrix clause is more ``canonically'' a
subject.
The effect is quite strong and the model is quite certain that when the matrix
verb encodes modality and the thematic role of `Agent' is assigned by the
infinitive verb, there is a big increase in the likelihood of attraction.

% section Thematic Role of the Controller (end)

\section{Dialect and Authorship} % (fold)

The model estimates that Attic and Jonic are quite different on how often they
employ case attraction, Jonic being predicted to use it much less often when all
the other factors are the same.
Although the model can not estimate a precise value (the standard deviation is
on the higher end), it is very confident of the contrast.
It also predicts that dialects should have a great amount of variation, as
denoted by the parameter $\sigma_\text{Dia}$.

\begin{center}
	
\begin{tabular}[t]{lrrrrr}
\toprule
Effect & mean & sd & 5.5\% & 94.5\% & rhat\\
\midrule
$\overline{\beta}_{\text{Dia}}$ & -0.37 & 0.80 & -1.64 & 0.93 & 1\\
$\sigma_{\text{Dia}}$ & 1.33 & 0.86 & 0.27 & 2.90 & 1\\
$\beta_{\text{Attic}}$ & 0.13 & 0.87 & -1.26 & 1.50 & 1\\
$\beta_{\text{Jonic}}$ & -1.60 & 1.14 & -3.50 & 0.17 & 1\\
$\beta_{\text{Attic}} - \beta_{\text{Jonic}}$ & 1.73 & 0.94 & -0.04 & 3.01 & --\\
\bottomrule
\end{tabular}

\end{center}

\begin{figure}[!h]
\centering
  \includegraphics[width=0.5\linewidth]{../analysis/figs/posterior_diff_dia.png}
	\caption{Posterior distribution of the difference between
		$\beta_\text{Attic}$ and $\beta_\text{Jonic}$ with the 89\%
	Highest Posterior Density Interval.}
\end{figure}
\clearpage

Finally, a word on authorship.
After controlling for dialect and preferred choice of matrix verb, most of the
authors are predicted to similar in their use of case attraction, with a curious
tendency of Xenophon to be, on average, less likely to employ it.

\begin{center}
	\includegraphics[width=0.75\linewidth]{../analysis/figs/precis_author.png}
\end{center}

% section Authorship and Dialect (end)

% chapter Analysis (end)

\renewcommand*{\mkbibnamefamily}[1]{\textsc{#1}}
\printbibliography%
\renewcommand*{\mkbibnamefamily}[1]{#1}

\end{document}
